\section{Ranh giới hệ thống, tác nhân và bên thụ hưởng}
\noindent \textbf{Ranh giới hệ thống.} Phạm vi phân tích bắt đầu khi người bệnh gửi thông tin triệu chứng qua ứng dụng, bao gồm quá trình AI hỏi thêm thông tin, tra cứu và soạn nội dung tư vấn, sau đó chuyển cho bác sĩ kiểm tra và phê duyệt. Quy trình kết thúc khi người bệnh nhận được tư vấn hoặc xác nhận lịch khám nếu cần. Các hoạt động khám trực tiếp, xét nghiệm, điều trị, thanh toán và cấp thuốc nằm ngoài phạm vi phân tích.

\noindent \textbf{Các tác nhân tham gia.} Người bệnh cung cấp thông tin triệu chứng và phản hồi các câu hỏi bổ sung; bác sĩ kiểm tra, điều chỉnh và phê duyệt nội dung tư vấn; nhân viên tiếp nhận hỗ trợ đặt lịch và xử lý các yêu cầu cần trợ giúp; bộ phận quản lý và kỹ thuật đảm bảo hệ thống vận hành ổn định. Ứng dụng tích hợp AI là công cụ hỗ trợ thu thập, xử lý và trao đổi thông tin giữa các bên.

\noindent \textbf{Bên thụ hưởng.} Hệ thống mang lại lợi ích cho người bệnh, bác sĩ, nhân viên y tế và cơ sở khám chữa bệnh. Cụ thể:
\begin{itemize}
    \item \textit{Người bệnh:} Có thể gửi câu hỏi và mô tả triệu chứng bất cứ lúc nào, nhận tư vấn đã được bác sĩ kiểm tra và dễ dàng đặt lịch khám
    \item \textit{Bác sĩ và nhân viên y tế:} Có sẵn bản tóm tắt triệu chứng, giảm thời gian hỏi đáp và tiết kiệm thời gian tư vấn
    \item \textit{Cơ sở y tế:} Dễ quản lý các yêu cầu tư vấn, sắp xếp lịch khám và phân công công việc cho nhân viên
\end{itemize}
